%=============================================================
%  第2章 先行研究
%=============================================================
\chapter{先行研究}
\label{chap:literature}

\section{企業ライフサイクル論}
\label{sec:lit-lifecycle}

\subsection{ライフサイクルの概念と測定}

企業ライフサイクル理論の源流は，\tcite{Mueller1972} の産業組織論的アプローチに遡る。
Mueller は，企業が誕生・成長・成熟・衰退という段階を経る過程で，
経営者の投資行動や利益分配が体系的に変化することを理論的に示した。
この枠組みは，企業の財務政策がライフサイクル段階に依存するという考え方の基礎を形成している。

\tcite{Anthony1992} は，会計情報の有用性がライフサイクル段階によって異なることを実証した。
同研究は売上高成長率・設備投資比率・配当支払率を代理変数として企業を成長期・成熟期・衰退期に
分類し，各段階で株価と利益・売上高の関連性が異なることを示した。
この知見は，企業の財務特性がライフサイクルに沿って構造的に変化することを支持する。

\tcite{Dickinson2011} は，営業・投資・財務キャッシュフローの符号パターンに基づいて
企業を5段階のライフサイクル（導入期・成長期・成熟期・動揺期・衰退期）に分類する
手法を提案した。同研究は，この分類が企業の収益性・成長性・配当政策と整合することを示し，
企業ライフサイクル代理変数として広く採用されるに至った。
本研究における5クラスター分類は，Dickinson の5段階分類と概念的に対応するが，
キャッシュフローの符号パターンではなく，11の財務変数に基づくデータ駆動型の分類である点で異なる。

\tcite{Habib2019survey} はライフサイクル研究を包括的にサーベイし，
会計・財務・コーポレートガバナンスの各分野でライフサイクル概念がどのように
応用されてきたかを整理した。
同サーベイは，ライフサイクル研究が配当政策，資本構造，リスクテイキング，M\&A，
資本コストなど多岐にわたる企業行動の説明に有効であることを確認している。

\subsection{財務政策とライフサイクル}

\paragraph{配当政策}
\tcite{DeAngelo2006} は，企業のライフサイクル段階が配当政策を規定するという
ライフサイクル仮説を実証的に検証した。同研究は，利益剰余金と払込資本の比率
（RE/TE ratio）が配当支払いの傾向を強く予測することを示し，
成熟企業ほど配当を支払う傾向があることを明らかにした。
\tcite{Fama2001} は米国上場企業において配当支払企業の割合が長期的に減少している現象を
記述し，その主因が企業特性の変化（小規模・低収益・高成長企業の増加）にあることを示した。
\tcite{Denis2008} は国際的なデータを用いて，配当支払いの決定要因がライフサイクル段階と
整合することを確認した。
これらの知見は，本研究のクラスター間で観察される配当関連指標の差異が，
ライフサイクル段階の差異を反映している可能性を示唆する。

\paragraph{資本構造}
ライフサイクルに沿った資本構造の変化も広く研究されている。
\tcite{Hirsch2011} は，企業のライフサイクル段階に応じて最適な負債水準が異なることを
理論的に示した。成長期の企業は内部資金と株式発行を優先し，
成熟期の企業はフリーキャッシュフローの規律づけとして負債を活用する傾向がある。
\tcite{Castro2016} は欧州上場企業のパネルデータを用いて，
目標レバレッジとその調整速度がライフサイクル段階によって異なることを実証した。
本研究において，\Cluster{2}（財務優良型）と \Cluster{0}（標準型）の自己資本比率の差異は，
こうしたライフサイクルに沿った資本構造の違いを反映していると考えられる。

\paragraph{リスクテイキング}
\tcite{Habib2017} は，企業のリスクテイキング行動がライフサイクル段階と投資家心理の
双方に依存することを示した。導入期・衰退期の企業はリスクテイキング水準が高く，
成熟期の企業はそれが低い傾向にある。
\tcite{Faff2016} は配当・投資・資金調達・現金保有の4つの財務政策が
ライフサイクルに沿って体系的に変化することを包括的に実証した。
これらの知見は，本研究のクラスター遷移が企業のリスク・プロファイルの変化と
関連している可能性を支持する。

\section{財務クラスタリングと機械学習}
\label{sec:lit-clustering}

\subsection{財務指標によるクラスタリング}

財務指標に基づく企業の分類は長い研究の歴史を持つ。
\tcite{Gupta1972} は財務比率のクラスター分析により産業特性を捉えることを試み，
\tcite{Pinches1973} は財務パターンの時系列的安定性を検証した。
これらの初期研究は，財務指標に基づく企業分類が経済的に意味のあるグループを
形成しうることを示している。

\tcite{GonzalezRuiz2021} は K-Means クラスタリングとマルコフ連鎖を組み合わせ，
企業の信用品質クラスター間の遷移確率を推定した。
本研究の方法論と最も類似しており，クラスタリング＋遷移分析の先駆的研究として位置づけられる。
ただし同研究は35社と小規模な分析に留まっており，日本の上場企業全体を対象とした
大規模分析は行われていない。

\tcite{Li2022} は PCA と K-Means の組み合わせにより企業財務リスクの識別を行い，
高リスク企業と低リスク企業のクラスターを識別することに成功した。

\tcite{Fodor2021} は株式リターンの相関構造に基づく財務クラスターが
産業分類とは異なるグループ構造を捉えることを示した。
この知見は，財務データに基づくクラスタリングが従来の産業分類では捉えきれない
企業間の類似性を明らかにしうることを示唆しており，
本研究が産業分類ではなくデータ駆動型のクラスタリングを採用する根拠の一つとなる。

\subsection{機械学習を用いた財務予測}

倒産予測の分野では，機械学習手法の優位性が実証されている。
\tcite{Shumway2001} はハザードモデルの枠組みを提案して以来，
動的モデルが静的モデルを上回ることが示されてきた。
近年では，勾配ブースティング手法が特に高い予測精度を達成している。

\tcite{Nguyen2023} は Random Forest, LightGBM, XGBoost, NGBoost を
比較し，XGBoost が AUC 約 94\% を達成することを示した。
同研究はさらに SHAP（Shapley Additive Explanations）を用いて
各変数の貢献度を定量化し，モデルの解釈可能性を確保した。

\tcite{BenJabeur2023} は特徴量重要度に基づく変数選択（FS-XGBoost）を提案し，
財務比率の選択がモデル精度に与える影響を分析した。
適切な特徴量選択がモデルの汎化性能を向上させることを示している。

\tcite{Zhang2023} は XGBoost に証拠理論を組み合わせた財務リスク早期警戒モデルを構築し，
SHAP により各変数の方向性を定量化した。

これらの研究は，XGBoost + SHAP の組み合わせが財務予測において
高精度かつ解釈可能なフレームワークを提供することを示している。
しかし，既存研究は倒産予測やリスク分類を対象としており，
財務クラスター間の遷移を予測する枠組みは提示されていない。

\section{本研究の位置づけ}
\label{sec:lit-position}

既存研究との対比を表\ref{tab:lit-comparison}に示す。

\begin{table}[htbp]
  \centering
  \caption{先行研究との比較}
  \label{tab:lit-comparison}
  \small
  \begin{tabularx}{\linewidth}{lXXcc}
    \toprule
    研究 & 手法 & データ & 遷移分析 & 要因解明 \\
    \midrule
    \tcite{Dickinson2011} & CF パターン & 米国上場 & なし & なし \\
    \tcite{GonzalezRuiz2021} & K-Means + Markov & 35社 & あり & なし \\
    \tcite{Li2022} & PCA + K-Means & 中国上場 & なし & なし \\
    \tcite{Nguyen2023} & XGBoost + SHAP & フランス企業 & なし & あり \\
    \textbf{本研究} & PCA + K-Means + XGBoost + SHAP & 日本上場 3,558社 & \textbf{あり} & \textbf{あり} \\
    \bottomrule
  \end{tabularx}
\end{table}

本研究は，以下の3点において既存研究を拡張する。
第一に，\num{3558} 社・15 年間の大規模な日本企業パネルデータを用いることで，
\tcite{GonzalezRuiz2021} の小規模分析（35社）を規模面で大幅に拡張する。
第二に，ハンガリアン法による年度間ラベル整合手法を導入し，
クラスターの時系列的な対応関係を客観的に確立する。
第三に，XGBoost + SHAP を財務クラスター遷移分析に適用し，
遷移の有無と遷移先の双方を予測・説明する統合的なフレームワークを構築する。
これにより，ライフサイクル理論で示唆される企業の段階的変化を，
データ駆動型の手法で実証的に捉えることを試みる。
